% Section 6.2, from lectures/L06/appendix/b_microcomputer.md.
%
% One correction to the source: the clock table in 6.2.8 said that 16 MHz is 16 000 000 edges a
% second, which is twice as many edges as there are; it is 16 000 000 periods, and the source is
% corrected to match. And 6.2.8 said the Uno's clock comes from a crystal; the ATmega328P on the
% Uno is clocked by a ceramic resonator (the crystal clocks the USB chip), as 11.1.7 says, and the
% source is corrected to match.

\section{Mikrodatorns uppbyggnad}
\label{c6:app:b}

\subsection{Vad en mikrodator är}
\label{c6:b1}
En \textbf{dator} är en maskin som utför en lista av instruktioner, ett \textbf{program}, och som
kan läsa och skriva data medan den gör det. Den består av en centralenhet som utför
instruktionerna, minne som håller programmet och datat, och in- och utenheter som förbinder den med
omvärlden.

En \textbf{mikrodator} är en dator där allt detta sitter i en eller några få kretsar. När hela
datorn, centralenhet, minnen och portar, sitter på \textbf{ett enda chip} kallas kretsen en
\textbf{mikrokontroller}. Den ATmega328P som sitter på ett Arduino Uno-kort, och som du
programmerar från kapitel~\ref{c7:ch}, är en mikrokontroller.

Mikrokontroller finns överallt där något ska styras: i tvättmaskiner, kaffebryggare, bilar, hissar,
fjärrkontroller och industrirobotar. En modern bil har hundratals. De kallas \textbf{inbyggda
system}, eftersom datorn är inbyggd i en produkt och sällan syns. Till skillnad från en persondator
kör en mikrokontroller oftast ett enda program, från det att strömmen slås på tills den slås av,
och programmet handlar om att läsa givare och styra saker: det som kursplanen kallar
\textbf{styrobjekt}.

\subsection{Blockschemat}
\label{c6:b2}
Alla datorer, från den minsta mikrokontroller till en server, är uppbyggda av samma block:

\bookfigure[tbp]{l06_microcomputer}{Mikrodatorns blockschema.}{c6:fig:microcomputer}

\begin{booktable}{@{}p{11.6em}L@{}}
  \thead{Block} & \thead{Uppgift} \\
  \midrule
  \textbf{CPU} (centralenhet) & Hämtar instruktioner och utför dem. Består av styrenhet, register
    och ALU. \\
  \textbf{Styrenhet} & Tolkar (avkodar) varje instruktion och styr de andra blocken så att den
    utförs. \\
  \textbf{Programräknare, PC} & Håller adressen till nästa instruktion. \\
  \textbf{Instruktionsregister, IR} & Håller den instruktion som utförs just nu. \\
  \textbf{Register} & Snabba lagringsplatser inuti CPU:n, där räkningarna görs. \\
  \textbf{ALU} & Aritmetisk-logisk enhet: adderar, subtraherar, jämför och gör logiska
    operationer. \\
  \textbf{SREG} & Statusregistret: flaggor som beskriver resultatet av den senaste räkningen. \\
  \textbf{Programminne} & Håller programmet, instruktion för instruktion. \\
  \textbf{Dataminne} & Håller variabler och annat data medan programmet kör. \\
  \textbf{In- och utportar} & Förbinder datorn med omvärlden: knappar och givare in, lysdioder
    och reläer ut. \\
  \textbf{Bussar} & Ledningarna mellan blocken, som adresser och data går på. \\
  \textbf{Klocka} & Ger takten som alla vippor i datorn uppdateras i. \\
\end{booktable}

Resten av det här avsnittet går igenom blocken ett i taget, och slutar med hur de ser ut i
ATmega328P.

\subsection{CPU:n: ALU, register och styrenhet}
\label{c6:b3}
\textbf{Registren} är små, snabba minnen inuti CPU:n. I ATmega328P finns 32 stycken, \code{r0}
till \code{r31}, med 8 bitar vardera: varje register är åtta D-vippor med gemensam klocka
(\secref{c6:a6}). All räkning sker i registren. Ett tal i dataminnet måste först hämtas in i ett
register innan något kan göras med det.

\textbf{ALU:n}, den aritmetisk-logiska enheten, är det enda ställe i datorn där något faktiskt
räknas. Den tar ett eller två tal från registren, utför en operation och lämnar resultatet i ett
register. Den är ett stort \textbf{kombinatoriskt nät}: samma sorts nät som i kapitel~\ref{c2:ch}
och \ref{c4:ch}, bara större. Kärnan i den är en \textbf{adderare}, byggd av heladderare
(\secref{c4:a8}):

\bookfigure{l06_full_adder}{En heladderare av två XOR-grindar, två AND-grindar och en
OR-grind.}{c6:fig:full_adder}

En heladderare adderar tre bitar: två siffror och minnessiffran från positionen till höger. Åtta
heladderare i rad, där varje \code{Cut} blir nästa positions \code{Cin}, adderar två 8-bitars tal,
precis som när du adderar två binära tal för hand i \secref{c1:a7}. Den sista minnessiffran, den
som inte får plats i 8 bitar, är \textbf{carry-flaggan}.

\textbf{Statusregistret, SREG}, är ett register där ALU:n sparar information om resultatet: blev
det noll, blev det en minnessiffra över, blev det negativt? De här bitarna kallas
\textbf{flaggor}. Programmet använder dem för att fatta beslut, \enquote{hoppa om resultatet blev
noll}, vilket är ämnet för kapitel~\ref{c8:ch} och \ref{c9:ch}.

\textbf{Styrenheten} läser instruktionen, avgör vad den betyder, och skickar
\textbf{styrsignaler} till de andra blocken: vilka register ALU:n ska läsa, vilken operation den
ska göra, vart resultatet ska. Den är i grunden ett sekvensnät, som stegar sig igenom varje
instruktion.

\subsection{Programräknaren och instruktionscykeln}
\label{c6:b4}
Ett program är en lista av instruktioner i programminnet, en efter en, var och en på sin adress.
\textbf{Programräknaren}, PC (\emph{program counter}), håller adressen till den instruktion som
står på tur. Den är en räknare (\secref{c6:a7}), och det är den som gör att datorn går framåt i
programmet.

Allt en CPU gör är att upprepa samma tre steg, \textbf{instruktionscykeln}:

\bookfigure{l06_instruction_cycle}{Instruktionscykeln: hämta, avkoda och
utför.}{c6:fig:instruction_cycle}

\begin{enumerate}
  \item \textbf{Hämta} (\emph{fetch}): läs instruktionen på den adress PC pekar på, lägg den i
    instruktionsregistret, och räkna upp PC så att den pekar på nästa instruktion.
  \item \textbf{Avkoda} (\emph{decode}): styrenheten tolkar bitarna i instruktionen. Vilken
    operation? Vilka register?
  \item \textbf{Utför} (\emph{execute}): ALU:n räknar, ett register skrivs, eller data flyttas
    till eller från minnet eller en port.
\end{enumerate}

Sedan börjar det om, med nästa instruktion. Ett litet exempel, tre instruktioner i
programminnet:

\begin{narrowtable}{@{}cll@{}}
  \thead{Adress} & \thead{Instruktion} & \thead{Betyder} \\
  \midrule
  0 & \code{ldi r16, 5} & lägg talet 5 i register r16 \\
  1 & \code{inc r16} & öka r16 med 1 \\
  2 & \code{rjmp 1} & hoppa till adress 1 \\
\end{narrowtable}

\begin{narrowtable}{@{}cclllc@{}}
  \thead{Steg} & \thead{PC före} & \thead{Hämtas} & \thead{Utförs} & \thead{r16 efter} &
    \thead{PC efter} \\
  \midrule
  1 & 0 & \code{ldi r16, 5} & r16 = 5 & 5 & 1 \\
  2 & 1 & \code{inc r16} & r16 = r16 + 1 & 6 & 2 \\
  3 & 2 & \code{rjmp 1} & PC = 1 & 6 & 1 \\
  4 & 1 & \code{inc r16} & r16 = r16 + 1 & 7 & 2 \\
  5 & 2 & \code{rjmp 1} & PC = 1 & 7 & 1 \\
\end{narrowtable}

Lägg märke till steg 3. Ett \textbf{hopp} är ingenting annat än en instruktion som skriver ett nytt
värde i programräknaren. Normalt räknar PC uppåt ett steg i taget; ett hopp låter programmet gå
tillbaka och upprepa något, eller hoppa över något. Alla loopar och alla beslut i ett program är
byggda på det, och det är vad kapitel~\ref{c9:ch} handlar om.

På ATmega328P tar de flesta instruktioner \textbf{en klockcykel}: hämtningen av nästa instruktion
sker samtidigt som den förra utförs. Vid 16~MHz blir det upp till sexton miljoner instruktioner i
sekunden.

\subsection{Minnen: programminne, dataminne och EEPROM}
\label{c6:b5}
En mikrodator har flera sorters minne, för olika ändamål:

\begin{narrowtable}{@{}lllc@{}}
  \thead{Minne} & \thead{Håller} & \thead{Bevaras utan ström?} & \thead{I ATmega328P} \\
  \midrule
  \textbf{Programminne} (flash) & programmet & ja & 32~kB \\
  \textbf{Dataminne} (SRAM) & variabler och stacken & nej & 2~kB \\
  \textbf{EEPROM} & inställningar som ska sparas & ja & 1~kB \\
\end{narrowtable}

\textbf{Flashminnet} håller programmet. Det skrivs när du programmerar kretsen, i
kapitel~\ref{c11:ch}, och finns kvar när strömmen slås av, så att programmet startar igen nästa
gång. Det kan skrivas om ungefär tiotusen gånger, vilket räcker gott för utveckling men inte för
data som ändras hela tiden.

\textbf{SRAM} (\emph{static RAM}) är dataminnet: variabler, mellanresultat och \textbf{stacken},
där datorn sparar returadresser när den anropar en subrutin (kapitel~\ref{c10:ch}). Det är snabbt
och kan skrivas hur många gånger som helst, men innehållet försvinner när strömmen bryts. Varje bit
i ett SRAM är i princip ett lås.

\textbf{EEPROM} bevarar sitt innehåll utan ström, som flashminnet, men kan skrivas en byte i taget.
Det används till sådant som ska överleva ett strömavbrott: en kalibrering, en inställning, en
räknare av driftstimmar. Boken använder det inte.

Lägg märke till att programminnet och dataminnet är \textbf{separata} minnen med var sin buss. Det
kallas en \textbf{Harvardarkitektur}, och det betyder att CPU:n kan hämta nästa instruktion
samtidigt som den läser eller skriver data. De flesta persondatorer har i stället program och data i
samma minne, en \textbf{von Neumann-arkitektur}. För dig som programmerar betyder skillnaden att ett
tal i programminnet och ett tal i dataminnet ligger på olika ställen och nås med olika
instruktioner, vilket du märker i kapitel~\ref{c7:ch}.

\subsection{Bussar}
\label{c6:b6}
Blocken är förbundna med \textbf{bussar}: grupper av ledningar som bär ett tal i taget, en bit per
ledning.
\begin{itemize}
  \item \textbf{Adressbussen} säger \textbf{var}: vilken minnesplats eller vilket register som ska
    läsas eller skrivas. Med \code{n} ledningar kan den peka ut \code{2^n} olika adresser.
    Programminnet i ATmega328P har 16\,384 instruktionsplatser, och $2^{14} = 16\,384$, så
    programräknaren behöver 14 bitar.
  \item \textbf{Databussen} bär \textbf{vad}: själva värdet som läses eller skrivs. I en 8-bitars
    dator som ATmega328P är den 8 bitar bred, och det är därför den kallas 8-bitars: den hanterar
    en byte i taget.
  \item \textbf{Styrsignalerna} säger \textbf{hur}: läs eller skriv, och när.
\end{itemize}

En buss är ingen ny sorts komponent. Den är ledningar, och i en mikrokontroller sitter de inuti
chipet. Men begreppet förklarar mycket: allt som ska flyttas mellan två block går över en buss, en
sak i taget, och det är därför en instruktion som flyttar data mellan minnet och ett register
ibland tar två klockcykler i stället för en.

\subsection{In- och utportar}
\label{c6:b7}
En dator som inte kan påverka något eller känna av något är ganska meningslös för den som ska styra
en maskin. \textbf{Portarna} är förbindelsen med omvärlden.

En \textbf{utport} är ett register vars utgångar är kopplade till kretsens \textbf{stift}. Skriver
programmet en etta i en bit, blir motsvarande stift 5~V; skriver det en nolla, blir stiftet 0~V.
Koppla en lysdiod med ett förkopplingsmotstånd till stiftet, precis som till en grindutgång i
Labb~2, så styr programmet lysdioden. Koppla ett relä via en transistor, så styr programmet en lampa
på 24~V eller en motor.

En \textbf{inport} är motsatsen: ett sätt för programmet att läsa av spänningen på ett antal stift,
som ettor och nollor. En knapp, en gränslägesbrytare eller en nivågivare blir då en bit som
programmet kan fatta beslut utifrån.

I ATmega328P kan varje stift vara antingen in- eller utgång, och det väljer programmet själv. Tre
register per port styr det: \code{DDRx} (riktningen), \code{PORTx} (vad en utgång ska ge) och
\code{PINx} (vad en ingång läser). Hur de används är ämnet för kapitel~\ref{c8:ch} och
\ref{c10:ch}.

\subsection{Klockan}
\label{c6:b8}
Allt i mikrodatorn går i takt med \textbf{klockan}, samma sorts fyrkantvåg som styr vipporna i
\secref{c6:a5}. Vid varje stigande flank tar varje register, programräknaren och styrenheten ett
steg.

På ett Arduino Uno-kort kommer klockan från en \textbf{keramisk resonator} på 16~MHz bredvid
chipet. Det ger:

\begin{narrowtable}{@{}ll@{}}
  \thead{Storhet} & \thead{Värde} \\
  \midrule
  Klockfrekvens & 16~MHz = 16\,000\,000 perioder per sekund \\
  En klockcykel & 1 / 16\,000\,000~s = 62,5~ns \\
  Instruktioner per sekund & upp till 16 miljoner (de flesta tar en cykel) \\
\end{narrowtable}

Att varje instruktion tar ett känt antal klockcykler gör att du kan räkna ut exakt hur lång tid en
del av ett program tar. Det är så en tidsfördröjning byggs, utan någon klocka på väggen: låt datorn
räkna ett visst antal varv, och räkna ut hur många cykler det blir. Det gör du i
kapitel~\ref{c9:ch}.

\subsection{ATmega328P och Arduino Uno}
\label{c6:b9}
Mikrokontrollern i resten av boken är \textbf{ATmega328P} från Microchip, samma krets som sitter på
ett \textbf{Arduino Uno}-kort. Den är en 8-bitars mikrokontroller i familjen \textbf{AVR}.

\bookfigure[tbp]{l06_atmega328p}{Förenklat blockschema för ATmega328P.}{c6:fig:atmega328p}

\begin{booktable}{@{}p{11.5em}L@{}}
  \thead{Egenskap} & \thead{ATmega328P} \\
  \midrule
  Arkitektur & 8-bitars AVR, Harvard \\
  Arbetsregister & 32 st, \code{r0}--\code{r31}, 8 bitar vardera \\
  Programminne (flash) & 32~kB, varav en liten del upptas av Arduinos bootloader \\
  Dataminne (SRAM) & 2~kB \\
  EEPROM & 1~kB \\
  Klocka på Arduino Uno & 16~MHz (keramisk resonator) \\
  Portar & B, C och D; 23 stift som kan vara in- eller utgångar \\
  Kringenheter & tre timrar, seriekommunikation (USART), analog-till-digital-omvandlare
    (ADC) \\
  Matning & 5~V på Arduino Uno \\
\end{booktable}

Blocken från \secref{c6:b2} finns alla där. CPU:n med sina 32 register, sin ALU och sitt SREG är
kärnan. Flashminnet håller programmet, SRAM variablerna. Portarna B, C och D är utportar och
inportar, och kortets numrerade stift (0--13, A0--A5) är kopplade till dem. Därtill kommer
\textbf{kringenheter}, egna små hårdvarublock för sådant som är vanligt i styrning: timrar som
räknar tid, en seriekanal för att prata med en dator, och en omvandlare som mäter analoga
spänningar. Boken använder portarna; resten får vänta till en fortsättningskurs.

Två saker är värda att veta innan du börjar programmera:
\begin{itemize}
  \item Arduino-kortets stiftnummer är \textbf{kortets} numrering, inte kretsens. Stift 13 på
    kortet, där den inbyggda lysdioden sitter, är bit~5 i port B, \code{PB5}. Kopplingen mellan
    dem kommer i kapitel~\ref{c8:ch}.
  \item Boken programmerar kretsen direkt i \textbf{assembler}, utan Arduinos egna
    programbibliotek. Varje instruktion du skriver blir en instruktion i flashminnet, och du ser
    exakt vad CPU:n gör.
\end{itemize}

\subsection{Från grind till dator}
\label{c6:b10}
Varje block i mikrodatorn går tillbaka till något du redan har byggt:

\begin{booktable}{@{}p{8.8em}Lp{9.8em}@{}}
  \thead{Block} & \thead{Byggt av} & \thead{Där du mötte det} \\
  \midrule
  ALU & kombinatoriska nät: adderare, AND, OR, XOR & kapitel~\ref{c2:ch} och \ref{c4:ch} \\
  Heladderare & XOR, AND, OR & \secref{c4:a8} \\
  Register & D-vippor med gemensam klocka & \secref{c6:a6} \\
  Programräknare & en räknare av vippor & \secref{c6:a7} \\
  Instruktionsregister & ett register & \secref{c6:a6} \\
  Statusregistret & ett register, en vippa per flagga & \secref{c6:a6} \\
  Dataminne (SRAM) & tusentals lås, och avkodare som väljer ut rätt byte & \secref{c6:a4},
    kapitel~\ref{c4:ch} \\
  Utport & ett register kopplat till stiften & \secref{c6:a6} \\
  Adressavkodning & kombinatoriska nät som väljer en adress & kapitel~\ref{c4:ch} \\
  Styrenhet & ett sekvensnät & \secref{c6:a1} \\
\end{booktable}

Ingenting i en dator är magiskt. Den är grindar och vippor, i stora mängder, i takt med en klocka.
Det du lär dig i resten av boken är hur man styr dem med ett program.
